\paragraph{The chain over every finite-type base.}
The smooth-base splitting above is enough to remove regularity from
the Kummer comparison itself.  One needs injectivity on sections,
not descent of a family of cohomological retractions along a
hypercover.

\begin{lemma}[Sections and their completion on a reduced base]
\label{lem:B-sections-reduced}
\status{Proved}
Let \(S/\F_q\) be reduced and of finite type, without a
separatedness assumption, and let \(A/S\) be an abelian scheme.
Put \(C_S=A(\bar S)\).  Then \(C_S/(C_S)_{\rm tors}\)
is finitely generated, its Galois action has finite image, and
the natural map gives an isomorphism
\[
 C_S\otimes_{\Z}\Q_\ell
 \xrightarrow{\sim}
 (C_S)^{\wedge}_\ell\otimes_{\Z_\ell}\Q_\ell
 \qquad (\ell\ne p).
\]
There exists a finite disjoint union \(V\) of smooth separated
finite-type \(\F_q\)-schemes with a jointly surjective map
\(v:V\to S\).  Pullback embeds \(C_S\) in \(C_V\)
and embeds their quotients by torsion.
\end{lemma}
\begin{proof}
Take a finite affine cover of \(S\).  Alter each reduced
irreducible component of each member of the cover.  De Jong's
Theorem~4.1 and Remark~4.2 give a smooth separated alteration,
since \(\F_q\) is perfect; their disjoint union is the
required \(V\)~\cite{c_dejong1996}.  The construction and
surjectivity persist after extension to \(\overline{\F}_q\).
If a section becomes zero on \(\bar V\), its equalizer with
the zero section is a closed subscheme of \(\bar S\) whose
underlying set is all of \(\bar S\).  Its ideal is therefore
contained in the nilradical and is zero.  Thus pullback is
injective.  Applying the same argument to an integer multiple of
a section shows injectivity modulo torsion as well.

For each connected component of \(\bar V\), restriction to
its generic point embeds its sections into the points of an
abelian variety over a finitely generated regular extension of
\(\overline{\F}_q\).  Lang--Neron, in the exact form of
Conrad's Theorems~2.1 and~7.1, says that the quotient by the
constant trace is finitely generated~\cite{BConrad2006}.  All
\(\overline{\F}_q\)-points of an abelian variety over
\(\overline{\F}_q\) are torsion: the variety and each point
descend to a finite field.  The quotient of its function-field
points by all torsion is consequently finitely generated.  A
subgroup of a finitely generated abelian group is finitely
generated.  Applying this twice, first to \(C_V\) and then to
\(C_S\), proves the assertion about sections modulo torsion.
Representatives of finitely many generators are defined over one
finite extension of \(\F_q\), which proves the finite-image
assertion.

Evaluation at one geometric point of each connected component of
\(\bar S\) injects \(C_S[\ell^\infty]\) into a finite
sum of groups \(A_{\bar s}[\ell^\infty]\).  Indeed a section
of a finite etale group scheme which vanishes at one point of a
connected scheme vanishes everywhere.  Thus the Pontryagin dual
of \(C_S[\ell^\infty]\) is a quotient of a finite free
\(\Z_\ell\)-module; it is finitely generated.  Equivalently,
this torsion group is a sum of finitely many copies of
\(\Q_\ell/\Z_\ell\) and a finite group.  Its
\(\ell\)-adic completion has zero rationalization.  The
prime-to-\(\ell\) torsion is \(\ell\)-divisible and has
zero completion.  Finally the extension of the torsion subgroup
by the finite free group \(C_S/(C_S)_{\rm tors}\) splits as
an extension of abelian groups.  Computing completions in that
splitting proves the displayed natural isomorphism; its statement
does not depend on the splitting chosen.
\end{proof}

\begin{theorem}[The finite-field chain for arbitrary finite-type schemes]
\label{thm:B-finite-chain-all-bases}
\status{Proved}
Let \(S/\F_q\) be any scheme of finite type, let \(A/S\)
be an abelian scheme, and let \(\ell\ne p\), including
\(\ell=2\) when \(p\ne2\).  No regularity, reducedness,
properness or separatedness of \(S\) is required.  Set
\[
 W_S=H^1_{\Bcts}(\bar S,V_\ell A),\qquad
 N_S=A(\bar S)\otimes\Q_\ell.
\]
The geometric Kummer sequence
\[
 0\longrightarrow N_S\longrightarrow W_S
 \longrightarrow V_\ell H^1_{\et}(\bar S,\bar A)
 \longrightarrow0
\]
has a \(\Gamma\)-equivariant splitting.  The natural map
\[
 H^1_{\et}(S,A)[\ell^\infty]\longrightarrow
 H^1_{\et}(\bar S,\bar A)[\ell^\infty]^\Gamma
\]
has finite kernel and cokernel.  The four conditions in
Proposition~\ref{prop:B-finite-chain-proper} are therefore
equivalent for every such \(S\).
\end{theorem}
\begin{proof}
First suppose that \(S\) is reduced and take \(V\) as in
Lemma~\ref{lem:B-sections-reduced}.  Kummer maps commute with
pullback.  Write \(i:N_S\hookrightarrow N_V\) for the
injection of that lemma and \(v^*:W_S\to W_V\) for
cohomological pullback.  The proof of
Proposition~\ref{prop:B-finite-chain-proper}, component by
component, supplies a \(\Gamma\)-equivariant retraction
\(r_V:W_V\to N_V\).  The action of \(\Gamma\) on
\(N_V\) factors through a finite group \(G\).  Choose any
linear retraction \(b:N_V\to N_S\) of \(i\) and average:
\[
 b_G=\frac1{|G|}\sum_{g\in G}g\,b\,g^{-1}.
\]
The scalar \(|G|^{-1}\) is available in \(\Q_\ell\),
even when \(\ell\mid |G|\).  Hence
\[
 r_S=b_G\,r_V\,v^*:W_S\longrightarrow N_S
 \quad\hbox{satisfies}\quad r_S|_{N_S}=1.
\]
This proves the splitting.  It does not require \(v^*\) to be
injective on the whole of \(W_S\), nor compatible retractions
on intersections of the covering alterations.

For a nonreduced \(S\), filter its nilpotent ideal by powers,
so that successive closed immersions have square-zero ideals
\(I\).  On the small etale site, smoothness of \(A\) gives
a surjection from the sheaf of its sections to the sheaf on the
thickening's reduction.  Its kernel for a square-zero step is
\(\operatorname{Lie}(A)\otimes I\), an additive sheaf killed
by \(p\).  The cohomology exact sequence and induction show
that the kernels and cokernels on \(H^0\) and \(H^1\) are
of bounded \(p\)-power exponent.  Such maps are isomorphisms
on \(\ell\)-primary torsion and after tensoring sections
with \(\Q_\ell\).  The finite etale sheaves
\(A[\ell^a]\) and their cohomology are unchanged by a
nilpotent thickening.  Thus the Kummer sequence and its splitting
reduce to those already proved for \(S_{\rm red}\).

It remains to verify the passage from geometric splitting to the
stated finite comparison.  Norms from finite constant extensions
give
\[
 A(S)\otimes\Q_\ell=N_S^\Gamma.
\]
Indeed a finite set of section representatives is defined over
one finite extension, and averaging over that extension descends
the invariant rational combination; injectivity is immediate.
Furthermore
\(H^0(\bar S,V_\ell A)_\Gamma=0\).  To see this,
evaluate at a finite set of geometric points, one on each
connected component, lying over closed points of \(S\).
This is injective on sections of a lisse sheaf.  After a finite
constant extension it is equivariant into a finite sum of Tate
modules of abelian varieties over finite fields.  Their
geometric Frobenius eigenvalues have weight \(-1\), so no
positive power has eigenvalue one.  Continuous
Hochschild--Serre therefore identifies
\[
 H^1_{\Bcts}(S,V_\ell A)=W_S^\Gamma.
\]
Taking invariants in the split geometric Kummer sequence and
comparing with global Kummer gives the natural isomorphism
\[
 V_\ell H^1_{\et}(S,A)
 \xrightarrow{\sim}
 \bigl(V_\ell H^1_{\et}(\bar S,\bar A)\bigr)^\Gamma.
\]
The group \(H^1(S,A)[\ell]\) is a quotient of the finite
group \(H^1(S,A[\ell])\), and the same holds over
\(\overline{\F}_q\).  Finite-coefficient etale finiteness
also holds for a nonseparated finite-type scheme by a finite
affine cover and its Cech spectral sequence.  Consequently both
\(\ell\)-primary groups are cofinitely generated.  Taking
invariants preserves this property, and the Tate module of the
invariants is the invariant Tate module, since invariants commute
with inverse limits.  A map between cofinitely generated
\(\ell\)-primary torsion groups with an isomorphism on
rational Tate modules has finite kernel and cokernel.  This proves
the comparison and all four equivalences.
\end{proof}

The transcription of 12:79 itself explicitly retains the condition
that the base be regular.  The relative-curve proposition above
therefore proves the requested Jacobian--Brauer equivalence in
that stated scope.  The present theorem removes the former
singular-base gap in the Kummer chain; it does not prove its
equivalent finiteness conditions for every nonconstant abelian
scheme.
